%% ============================================================
%% CHAPTER 7 — RQ4: OFF-DESIGN
%% ============================================================
\chapter{Off-Design and Part-Load Performance (RQ4)}
\label{ch:results_rq4}

This chapter answers RQ4. Section~\ref{sec:rq4_partload} reports efficiency
retention at five load levels for each cycle at each site design point, and
Section~\ref{sec:rq4_flexibility} the Marchionni flexibility scores.

\section{Part-Load Efficiency Curves}
\label{sec:rq4_partload}

\begin{figure}[H]
\centering
\thesisfig{RQ4_OFFDESIGN_PARTLOAD.png}{1.0}{10cm}
\caption{Efficiency retention versus load at the three site design points
(Stodola correction with sliding-pressure control) and the Marchionni
flexibility scores.}
\label{fig:offdesign}
\end{figure}

\begin{table}[H]
\centering
\caption{Thermal efficiency (\%) at part load, and retention relative to the
design point (in parentheses).}
\label{tab:partload}
\renewcommand{\arraystretch}{1.25}
\footnotesize
\begin{tabularx}{\textwidth}{>{\raggedright\arraybackslash}p{0.9cm}>{\raggedright\arraybackslash}p{1.0cm}
  >{\centering\arraybackslash}X >{\centering\arraybackslash}X >{\centering\arraybackslash}X >{\centering\arraybackslash}X >{\centering\arraybackslash}X}
\toprule
\rowcolor{tableheader}
\textbf{Site} & \textbf{Cycle} & \textbf{100\,\%} & \textbf{80\,\%} & \textbf{60\,\%} & \textbf{40\,\%} & \textbf{20\,\%} \\
\midrule
\rowcolor{rowgray} CS2 & SRC     & 14.44 & 13.24 (0.917) & 11.49 (0.796) & 9.09 (0.629)  & 5.55 (0.384) \\
                   CS2 & sCO$_2$ & 11.43 & 11.50 (1.006) & 10.33 (0.904) & 8.00 (0.700)  & 4.68 (0.410) \\
\rowcolor{rowgray} CS2 & tCO$_2$ & 15.67 & 14.70 (0.938) & 13.18 (0.841) & 11.17 (0.713) & 8.75 (0.558) \\
\midrule
                   CS1 & SRC     & 23.67 & 22.49 (0.950) & 20.55 (0.868) & 17.76 (0.751) & 13.76 (0.582) \\
\rowcolor{rowgray} CS1 & sCO$_2$ & 22.24 & 21.09 (0.948) & 18.66 (0.839) & 14.99 (0.674) & 10.13 (0.455) \\
                   CS1 & tCO$_2$ & 24.97 & 23.95 (0.959) & 22.21 (0.889) & 19.79 (0.792) & 16.73 (0.670) \\
\midrule
\rowcolor{rowgray} CS3 & SRC     & 30.82 & 29.68 (0.963) & 27.63 (0.897) & 24.59 (0.798) & 20.21 (0.656) \\
                   CS3 & sCO$_2$ & 27.58 & 25.89 (0.939) & 22.91 (0.831) & 18.61 (0.675) & 12.99 (0.471) \\
\rowcolor{rowgray} CS3 & tCO$_2$ & 28.49 & 27.21 (0.955) & 25.18 (0.884) & 22.42 (0.787) & 18.97 (0.666) \\
\bottomrule
\end{tabularx}
\end{table}

\begin{findingbox}
\textbf{Finding 28 --- The tCO$_2$ cycle retains the most efficiency at deep
part load; the sCO$_2$ cycle loses the most.}
At 20\,\% load tCO$_2$ retains \keyresult{56--67\,\%} of its design
efficiency, steam 38--66\,\% and sCO$_2$ 41--47\,\%. The tCO$_2$ advantage
comes from the liquid pump, whose work falls in proportion to flow while
the Stodola penalty applies only to the turbine, and from the gentler
sliding-pressure schedule ($0.6 + 0.4f$) that a pumped cycle can tolerate.
The sCO$_2$ cycle carries the Stodola penalty on both its turbine and its
compressor, and because the compressor already absorbs 31--64\,\% of
turbine work at design, a proportionate loss in both machines erodes the net
output fastest.
\end{findingbox}

\begin{findingbox}
\textbf{Finding 29 --- The steam cycle's part-load behaviour depends strongly
on its design pressure: it is the most robust cycle at 500\,\textdegree C and
the least at 200\,\textdegree C.}
At CS3, where the boiler runs at 3.7\,MPa, steam retains 90\,\% at 60\,\%
load and 66\,\% at 20\,\% load, marginally better than tCO$_2$. At CS2 the
pinch-limited boiler pressure is only 0.16\,MPa, and sliding it to
0.03\,MPa at 20\,\% load leaves little expansion ratio above the 10\,kPa
condenser: retention falls to \keyresult{38\,\%}. Low-grade steam WHR is
therefore poorly suited to a fluctuating source.
\end{findingbox}

\begin{findingbox}
\textbf{Finding 30 --- The apparent efficiency \emph{gain} of sCO$_2$ at
80\,\% load at the 200\,\textdegree C site is a pressure-ratio effect, not a
near-critical effect.}
At CS2 the sCO$_2$ retention at 80\,\% load is 1.006. The sliding-pressure
schedule lowers the high-side pressure from 20 to 18\,MPa, and at a
190\,\textdegree C turbine inlet the design pressure ratio of 2.5 is above
the efficiency-optimal ratio, so the reduction helps more than the Stodola
penalty hurts. The compressor inlet state is fixed at 36\,\textdegree C and
8\,MPa at all loads, so the Interim Report's attribution of this feature to
near-critical compressor behaviour is withdrawn. The effect vanishes at
350 and 500\,\textdegree C, where the design pressure ratio is closer to
optimal.
\end{findingbox}

\section{Operational Flexibility Scoring}
\label{sec:rq4_flexibility}

\begin{table}[H]
\centering
\caption{Marchionni flexibility scores (retention values averaged over the
three sites).}
\label{tab:flexibility}
\renewcommand{\arraystretch}{1.3}
\begin{tabularx}{\textwidth}{>{\raggedright\arraybackslash}p{1.5cm}
  >{\centering\arraybackslash}X >{\centering\arraybackslash}X >{\centering\arraybackslash}X
  >{\centering\arraybackslash}X >{\centering\arraybackslash}X}
\toprule
\rowcolor{tableheader}
\textbf{Cycle} & \textbf{Turndown (0.30)} & \textbf{$\eta$ at 60\,\% (0.30)} &
\textbf{$\eta$ at 40\,\% (0.25)} & \textbf{Start-up (0.15)} & \textbf{Total} \\
\midrule
\rowcolor{rowgray}
SRC     & 0.70 & 0.854 & 0.726 & 0.55 & 0.730 \\
sCO$_2$ & 0.55 & 0.858 & 0.683 & 0.75 & 0.706 \\
\rowcolor{rowgray}
tCO$_2$ & 0.70 & 0.871 & 0.764 & 0.80 & \textbf{0.782} \\
\bottomrule
\end{tabularx}
\end{table}

The tCO$_2$ cycle scores highest (0.782), steam second (0.730) and sCO$_2$
third (0.706). The ordering is the same as the Interim Report's, but the
margins are smaller (the interim scores were 0.854, 0.763 and 0.801 for
tCO$_2$, steam and sCO$_2$) because the corrected recuperator reduces the
part-load advantage of both CO$_2$ cycles and the pinch-limited steam
pressures change the steam curves. The tCO$_2$ lead rests on three
properties: a liquid pump, the simplest start-up of the three (no steam drum,
no near-critical compressor control), and the best efficiency retention at
60 and 40\,\% load.

\section{Summary of RQ4 Results}
\label{sec:rq4_summary}

For a fluctuating heat source the transcritical CO$_2$ cycle is the most
flexible of the three, retaining 56--67\,\% of its design efficiency at
one-fifth load and scoring 0.78 on the Marchionni scale against 0.73 for
steam and 0.71 for sCO$_2$. The steam cycle is nearly as flexible at high
source temperature but degrades badly at 200\,\textdegree C, where its boiler
pressure is already very low. The sCO$_2$ cycle degrades fastest at deep
part load at every site. Flexibility is the one criterion on which tCO$_2$
is unambiguously preferred, and it is most valuable precisely at the
batch-type sources (electric arc furnaces, intermittent process trains)
where it matters.
